\section*{Bab \ref{ch:b1:det} --- Determinan dan Sistem Linear}

\begin{solution}{exo:b1:det:1}
$3 \times 2 - 1 \times 5 = 1$.

Yang kedua: $L_2 \leftarrow L_2 - L_1$, $L_3 \leftarrow L_3 - L_2$ (pada
baris aslinya) memberikan baris $(1,2,3), (3,3,3), (3,3,3)$: yaitu dua baris
yang sama, jadi determinannya $0$. (Sarrus menegaskan: $45 + 84 + 96 - 105 - 48 -
72 = 0$.)

Yang ketiga: ia Vandermonde dengan $x = 1, 2, 3$
(\cref{ex:b1:det:vandermonde}): $(2-1)(3-1)(3-2) = 2$.
\end{solution}

\begin{solution}{exo:b1:det:2}
Determinannya sama dengan (tambahkanlah semua kolomnya pada yang pertama, lalu faktorkanlah)
$(\lambda + 2)$ kali
\[
\begin{vmatrix}
1 & 1 & \lambda\\ 1 & \lambda & 1\\ 1 & 1 & 1
\end{vmatrix}
= -(\lambda - 1)^2
\]
(bersihkanlah dengan $L_1 \leftarrow L_1 - L_3$, $L_2 \leftarrow L_2 - L_3$
lalu jabarkanlah), yang memberikan $\det = -(\lambda+2)(\lambda-1)^2$. \omterm{def:b1:vspaces:free}{Basis}
$\iff \det \neq 0 \iff \lambda \notin \{1, -2\}$.
\end{solution}

\begin{solution}{exo:b1:det:3}
Sistem pertama: $\det = -7$;
$x = \frac{1}{-7}\begin{vmatrix} 5 & 1\\ 4 & -2\end{vmatrix}
= \frac{-14}{-7} = 2$,
$y = \frac{1}{-7}\begin{vmatrix} 2 & 5\\ 3 & 4\end{vmatrix}
= \frac{-7}{-7} = 1$. Periksa: $2(2) + 1 = 5$; $3(2) - 2 = 4$.

Sistem kedua: setelah $L_2 - L_1$ dan $L_3 - 2L_1$, barisnya menjadi
$(1,1,1)$, $(0,-2,0)$, $(0,-1,-3)$, sehingga
\[
\det A = \begin{vmatrix} 1&1&1\\ 1&-1&1\\ 2&1&-1\end{vmatrix}
= 1 \times \begin{vmatrix} -2 & 0\\ -1 & -3\end{vmatrix} = 6 .
\]
Cramer, dengan mengganti kolomnya oleh $(6,2,1)^{\mathsf T}$:
\[
x = \frac{6}{6} = 1, \qquad
y = \frac{12}{6} = 2, \qquad
z = \frac{18}{6} = 3
\]
(dengan pembilangnya dihitung dengan cara yang sama). Periksa: $1 + 2 + 3 = 6$; $1 -
2 + 3 = 2$; $2 + 2 - 3 = 1$.
\end{solution}

\begin{solution}{exo:b1:det:4}
Reduksikanlah matriks yang diperbesar: $L_2 \leftarrow L_2 - 2L_1$, $L_3
\leftarrow L_3 - L_1$:
\[
\begin{pmatrix}
1 & 2 & -1 & 1 & 1\\
0 & 0 & 3 & -3 & 3\\
0 & 0 & 3 & -3 & 3
\end{pmatrix}
\to
\begin{pmatrix}
1 & 2 & -1 & 1 & 1\\
0 & 0 & 1 & -1 & 1\\
0 & 0 & 0 & 0 & 0
\end{pmatrix}.
\]
Variabel porosnya $x, z$; variabel bebasnya $y, t$. Penyulihan mundur: $z =
1 + t$, $x = 1 - 2y + z - t = 2 - 2y$. \omterm{def:b1:logic:sets}{Himpunan} penyelesaiannya:
\[
\{(2 - 2y,\; y,\; 1 + t,\; t) : y, t \in \R\}
= (2, 0, 1, 0) + \operatorname{Vect}\bigl((-2,1,0,0),\,
(0,0,1,1)\bigr),
\]
yaitu bidang afin (berdimensi $2 = 4 - \operatorname{rk} 2$) pada $\R^4$.
\end{solution}

\begin{solution}{exo:b1:det:5}
Induksi; $n = 1$ berupa hasil kali kosong $= 1$. Untuk langkahnya, jalankanlah
$C_k \leftarrow C_k - x_1 C_{k-1}$ bagi $k = n, n-1, \dots, 2$ (dalam
urutan ini, sehingga setiap operasinya memakai kolom yang belum terubah). Adapun
baris pertamanya menjadi $(1, 0, \dots, 0)$; sedangkan pada baris $i \geq 2$, entri
ke-$k$-nya menjadi $x_i^{k-1} - x_1 x_i^{k-2} = x_i^{k-2}(x_i - x_1)$.
Dengan menjabarkannya sepanjang baris pertamanya lalu memfaktorkan $(x_i - x_1)$ keluar dari
setiap baris $i$:
\[
V(x_1, \dots, x_n)
= \prod_{i=2}^{n} (x_i - x_1)\cdot V(x_2, \dots, x_n),
\]
dan hipotesis induksinya melengkapkan hasil kali $\prod_{i<j}(x_j
- x_i)$.
\end{solution}

\begin{solution}{exo:b1:det:6}
Menjabarkan $D_n$ sepanjang baris pertamanya: $D_n = 2 D_{n-1} -
1\cdot\begin{vmatrix} 1 & \ast\\ 0 & D_{n-2}\text{-blok}
\end{vmatrix}$; dan determinan keduanya, yang dijabarkan sepanjang kolom
pertamanya, adalah $D_{n-2}$. Jadi $D_n = 2D_{n-1} - D_{n-2}$, yaitu $D_n -
D_{n-1} = D_{n-1} - D_{n-2}$: sehingga selisihnya tetap, sama dengan
$D_2 - D_1 = 1$. Maka $D_n = D_1 + (n - 1) = n + 1$.
(Periksa: $D_2 = 3$, dan kasus $3\times3$-nya adalah
\cref{ex:b1:det:cofactor}: $D_3 = 4$.)
\end{solution}

\begin{solution}{exo:b1:det:7}
$C_1 \leftarrow C_1 + C_2 + C_3$ membuat kolom pertamanya tetap
$(m+2)$; faktorkanlah ia:
\[
\det = (m+2)\begin{vmatrix}
1 & 1 & m\\ 1 & m & 1\\ 1 & 1 & 1
\end{vmatrix}
\overset{L_1 - L_3,\ L_2 - L_3}{=}
(m+2)\begin{vmatrix}
0 & 0 & m-1\\ 0 & m-1 & 0\\ 1 & 1 & 1
\end{vmatrix}
= (m+2)\cdot\bigl(-(m-1)^2\bigr)
\]
(jabarkanlah sepanjang kolom pertamanya: entri tunggalnya $1$ bertanda
$+$, dan determinan $2 \times 2$ sisanya adalah $0 \cdot 0 -
(m-1)(m-1) = -(m-1)^2$).

Pembahasannya. $m \notin \{1, -2\}$: yaitu sistem Cramer; dan lewat kesetangkupan
persamaannya, $x = y = z$, sehingga setiap persamaannya memberikan $(m + 2)x = 1$:
jadi penyelesaian tunggal $\bigl(\frac{1}{m+2}, \frac{1}{m+2},
\frac{1}{m+2}\bigr)$. $m = 1$: ketiga persamaannya semua terbaca $x + y
+ z = 1$: sehingga penyelesaiannya membentuk bidang afin $x + y + z = 1$. $m = -2$:
menjumlahkan ketiga persamaannya memberikan $0 = 3$: jadi \omterm{def:b1:logic:sets}{himpunan} penyelesaiannya kosong.
\end{solution}

\begin{solution}{exo:b1:det:8}
($\Rightarrow$) Jika $A^{-1}$ berentri bilangan bulat: maka $\det A \cdot \det
A^{-1} = 1$ dengan kedua determinannya bulat (karena jumlah hasil kali
entrinya): dan dua bilangan bulat yang hasil kalinya $1$ keduanya $\pm1$.

($\Leftarrow$) Rumus kofaktornya $A^{-1} = \frac{1}{\det
A}\operatorname{Com}(A)^{\mathsf T}$ (yang diperiksa untuk $n \leq 3$ lewat
penjabaran langsung, dan diakui pada umumnya) mempunyai
$\operatorname{Com}(A)$ yang berentri bulat (karena setiap kofaktornya
determinan yang bulat); dan \omterm{def:b1:arith:divides}{membaginya} dengan $\det A = \pm 1$ menjaga kebulatannya.
\end{solution}

\begin{solution}{exo:b1:det:9}
Tambahkanlah semua kolomnya pada yang pertama: maka setiap entri kolom pertamanya yang baru adalah
$a + nb$; faktorkanlah ia keluar, sehingga kolom pertamanya semuanya satu. Lalu
\omterm{met:b1:matrices:gauss}{operasi baris} $L_i \leftarrow L_i - L_1$ ($i \geq 2$) membersihkan
setiap entri di bawah $1$ kiri atasnya dan menyisakan $a$ pada diagonalnya dan
$0$ di tempat lain pada baris itu: sehingga matriksnya segitiga atas dengan
diagonal $(1, a, \dots, a)$. Maka
\[
\det(aI + bJ) = (a + nb)\, a^{\,n-1} ,
\]
yang taknol jika dan hanya jika $a \neq 0$ dan $a + nb \neq 0$: yaitu syarat pada
\cref{exo:b1:matrices:9}.
\end{solution}

\begin{solution}{exo:b1:det:10}
Operasi blok (yang masing-masingnya komposisi atas $n$ operasi skalar
yang bersesuaian, dan diperkenankan oleh \cref{thm:b1:det:props}~(1)):
\[
\begin{vmatrix} A & B\\ B & A\end{vmatrix}
\overset{C_1 \leftarrow C_1 + C_2}{=}
\begin{vmatrix} A + B & B\\ A + B & A\end{vmatrix}
\overset{L_2 \leftarrow L_2 - L_1}{=}
\begin{vmatrix} A + B & B\\ 0 & A - B\end{vmatrix}
= \det(A+B)\,\det(A-B),
\]
dengan memakai kaidah blok segitiganya pada langkah terakhirnya.
\end{solution}

\begin{solution}{exo:b1:det:11}
$C_1 \leftarrow C_1 + C_2 + C_3$ membuat kolom pertamanya tetap
$(a + b + c)$; faktorkanlah ia keluar, lalu $L_2 \leftarrow L_2 - L_1$,
$L_3 \leftarrow L_3 - L_1$:
\[
\Delta = (a+b+c)\begin{vmatrix}
1 & b & c\\
0 & a - b & b - c\\
0 & c - b & a - c
\end{vmatrix}
= (a+b+c)\bigl[(a-b)(a-c) + (b-c)^2\bigr],
\]
dan dengan menjabarkannya, $(a-b)(a-c) + (b-c)^2 = a^2 + b^2 + c^2 - ab - bc
- ca$. Atas $\C$, dengan $j^3 = 1$ dan $1 + j + j^2 = 0$:
\begin{align*}
(a + jb + j^2c)(a + j^2b + jc)
&= a^2 + b^2 + c^2 + (j + j^2)(ab + bc + ca)\\
&= a^2 + b^2 + c^2 - ab - bc - ca ,
\end{align*}
dan dari situlah pemfaktoran lengkapnya. (Secara struktur: kolom
$(1, j, j^2)^{\mathsf T}$ memenuhi $M\,(1, j, j^2)^{\mathsf T}
= (a + jb + j^2c)(1, j, j^2)^{\mathsf T}$, dan serupa bagi
$j^2$ dan $1$: sehingga ketiga faktornya adalah ketiga ``nilai eigen''
sirkulannya, yaitu kisah yang disistematiskan pada jilid Tahun~2.)
\end{solution}

\begin{solution}{exo:b1:det:12}
Tulislah $r = \operatorname{rk} A$.

\emph{Ada submatriks $r \times r$ yang dapat dibalik.} Pilihlah $r$
kolom \omterm{def:b1:vspaces:free}{bebas} pada $A$ lalu misalkan $B \in \mathcal{M}_{n,r}$
matriks yang dibentuknya: maka $\operatorname{rk} B = r$. Dan karena rank barisnya
sama dengan rank kolomnya (\cref{thm:b1:matrices:rank}), $B$ mempunyai $r$
baris \omterm{def:b1:vspaces:free}{bebas}; sehingga menyimpan baris itu menghasilkan submatriks $r \times r$
pada $A$ yang ber-rank $r$, yaitu dapat dibalik.

\emph{Tak ada yang lebih besar.} Misalkan $S$ submatriks $s \times s$
yang dapat dibalik, yang diambil dari kolom $j_1, \dots, j_s$ dan baris $i_1,
\dots, i_s$ pada $A$. Jika sebuah kombinasi $\sum_k \lambda_k
C_{j_k} = 0$ atas kolom \emph{penuh} yang bersesuaian lenyap,
maka membaca hanya baris $i_1, \dots, i_s$ memberikan $\sum_k
\lambda_k S_k = 0$ pada kolom $S$, sehingga semua $\lambda_k = 0$
(karena $S$ dapat dibalik): jadi kolom $C_{j_1}, \dots, C_{j_s}$ pada $A$
\omterm{def:b1:vspaces:free}{bebas}, dan $s \leq \operatorname{rk} A = r$.

Maka $\operatorname{rk} A$ tepat merupakan ukuran terbesar sebuah
submatriks yang dapat dibalik.
\end{solution}

\begin{solution}{pb:b1:det:1}
\textbf{1.} $V(1,2,3,4) = (2-1)(3-1)(4-1)(3-2)(4-2)(4-3) = 1
\cdot 2\cdot 3\cdot 1\cdot 2\cdot 1 = 12$. Adapun interpolasi pada
simpul yang berbeda meminta koefisien $P$ yang memecahkan $W c =
y$ dengan $W = (x_i^{\,j-1})$, dan $\det W = V \neq 0$: jadi sistem
Cramer.

\textbf{2.} Kerjakanlah kolomnya dari kiri ke kanan. $C_1$
adalah kolom tetap $P_0(x_i) = 1$ (dengan $P_0$ \omterm{def:b1:poly:def}{monik} berderajat
$0$). Andaikan kolom $1, \dots, j-1$ sudah tereduksi
menjadi pangkat murni $1, x_i, \dots, x_i^{\,j-2}$. Karena $P_{j-1}
= X^{j-1} + \sum_{k < j-1}\alpha_k X^k$, mengurangkan dari $C_j$
kombinasi $\sum_k \alpha_k\,(\text{kolom atas } x_i^k)$
--- yaitu operasi yang tak mengubah determinannya --- menyisakan
kolom pangkat murninya $x_i^{\,j-1}$. Setelah kolom terakhirnya
matriksnya menjadi matriks Vandermonde: jadi $\det = V(x_1, \dots, x_n)$.

\textbf{3.} \omterm{def:b1:poly:def}{Polinomial} $(j-1)!\,B_{j-1}$ bersifat \omterm{def:b1:poly:def}{monik} dan
berderajat $j - 1$, sehingga pertanyaan 2 memberikan
\[
\det\bigl(B_{j-1}(m_i)\bigr)
= \frac{V(m_1, \dots, m_n)}{0!\,1!\cdots(n-1)!} .
\]
Ruas kirinya adalah determinan sebuah matriks berentri \emph{bilangan bulat}
(karena $B_k$ bernilai bulat pada $\Z$: yaitu pertanyaan 16--17 pada
soal akhir pekan \cref{pb:b1:vspaces:1}), sehingga ia bilangan bulat;
dan ia positif karena $V(m_1, \dots, m_n) > 0$ untuk $m_1 <
\dots < m_n$. Jadi superfaktorialnya \omterm{def:b1:arith:divides}{membagi} hasil kali semua
selisih berpasangannya.

\textbf{4.} Faktorkanlah $x_i$ keluar dari setiap baris $i$:
$\det(x_i^{\,j})_{j = 1..n} = x_1\cdots x_n\,
\det(x_i^{\,j-1}) = x_1\cdots x_n\,V$.

\textbf{5.} $(W^{\mathsf T}W)_{ij} = \sum_k x_k^{\,i-1}
x_k^{\,j-1} = p_{i+j-2}$: jadi $S = W^{\mathsf T}W$. Maka $\det S =
\det(W^{\mathsf T})\det W = V^2$
(\cref{thm:b1:det:props}~(2),(4)). Untuk $x_i$ yang real: $\det S = V^2
\geq 0$, dan $S$ dapat dibalik jika dan hanya jika $V \neq 0$ jika dan hanya jika $x_i$
berbeda berpasangan --- yaitu uji bertanda tegas yang dapat dihitung dari
jumlah pangkatnya belaka.

\textbf{6.} Syarat interpolasinya $\sum_{k}
c_k\,x_i^{\,k} = y_i$ membentuk sistem $Wc = y$; dan $\det W = V \neq
0$ memberikan keberadaan dan ketunggalannya sekaligus. Inilah bukti
ketiga di dalam buku ini: yaitu rumus eksplisit pada
\cref{thm:b1:poly:lagrange}, argumen kernel pada
\cref{ex:b1:linmaps:interpolation}, dan Cramer di sini.

\textbf{7.} Cramer: $c_{n-1} = \det W'/\det W$ dengan $W'$ berupa $W$
yang kolom terakhirnya diganti oleh $y$. Menjabarkan $\det W'$ sepanjang
kolom itu:
\[
\det W' = \sum_{i=1}^n (-1)^{i+n} y_i\,V(x_1, \dots, \widehat{x_i},
\dots, x_n) .
\]
Kini $V = V(\setminus i)\cdot\prod_{j<i}(x_i - x_j)\prod_{j>i}
(x_j - x_i)$, dan mengalihkan hasil kali keduanya berongkos
$(-1)^{n-i}$:
\[
(-1)^{i+n}\,\frac{V(\setminus i)}{V}
= \frac{(-1)^{i+n}(-1)^{n-i}}{\prod_{j\neq i}(x_i - x_j)}
= \frac{1}{\prod_{j\neq i}(x_i - x_j)} ,
\]
dan dari situlah $c_{n-1} = \sum_i y_i/\prod_{j \neq i}(x_i - x_j)$ ---
yaitu rumus \omterm{pb:b1:vspaces:1}{selisih terbagi} lagi.

\textbf{8.} $L_3 \leftarrow L_3 - L_1$ memberikan baris $(1, x_1,
x_1^2)$, $(0, 1, 2x_1)$, $(0,\ x_2 - x_1,\ (x_2-x_1)(x_2+x_1))$;
lalu menjabarkannya sepanjang kolom pertamanya dan memfaktorkan $(x_2 - x_1)$:
\[
(x_2 - x_1)\begin{vmatrix} 1 & 2x_1\\ 1 & x_2 + x_1
\end{vmatrix} = (x_2 - x_1)(x_2 - x_1) = (x_2 - x_1)^2 .
\]
Ia taknol untuk $x_1 \neq x_2$: sehingga \omterm{def:b1:det:system}{sistem linear} yang menyatakan
$P(x_1) = u$, $P'(x_1) = v$, $P(x_2) = w$ pada koefisien
$P \in \R_2[X]$ bersifat Cramer --- jadi interpolasi Hermite dengan
simpul yang berganda terajukan dengan baik.

\textbf{9.} $P = a + bX + cX^2$ dengan $a = P(0) = 1$, $b = P'(0)
= 0$, $a + b + c = P(1) = 2$: jadi $c = 1$, sehingga $P = 1 + X^2$, yang tunggal.
Keselarasannya: di sini $x_1 = 0$, $x_2 = 1$ dan determinan
pertanyaan 8 adalah $(1 - 0)^2 = 1 \neq 0$.

\textbf{10.} Perhitungan langsungnya:
\[
C_2 = \frac{1}{(a_1+b_1)(a_2+b_2)} - \frac{1}{(a_1+b_2)(a_2+b_1)}
= \frac{(a_1+b_2)(a_2+b_1) - (a_1+b_1)(a_2+b_2)}
{\prod_{i,j}(a_i+b_j)} ,
\]
dan pembilangnya menjabar menjadi $a_1b_1 + a_2b_2 - a_1b_2 - a_2b_1
= (a_2 - a_1)(b_2 - b_1)$: yaitu selisih dibagi jumlah.

\textbf{11.} Untuk $i < n$, entri baru pada baris $i$ adalah
\[
\frac{1}{a_i + b_j} - \frac{1}{a_n + b_j}
= \frac{a_n - a_i}{(a_i + b_j)(a_n + b_j)} .
\]
Faktorkanlah $(a_n - a_i)$ keluar dari setiap baris $i < n$, lalu
$\frac1{a_n + b_j}$ keluar dari setiap kolom $j$: maka yang tersisa berentri
$\frac1{a_i + b_j}$ pada baris $i < n$ dan tetap $1$ pada
baris $n$ --- yaitu matriks $M$, dengan prafaktor yang diumumkan.

\textbf{12.} Pada $M$, untuk $j < n$ operasi $C_j \leftarrow
C_j - C_n$ mengubah baris $n$ menjadi $(0, \dots, 0, 1)$ dan, pada baris $i
< n$,
\[
\frac{1}{a_i + b_j} - \frac{1}{a_i + b_n}
= \frac{b_n - b_j}{(a_i + b_j)(a_i + b_n)} .
\]
Faktorkanlah $(b_n - b_j)$ dari setiap kolom $j < n$ dan $\frac1{a_i +
b_n}$ dari setiap baris $i < n$, lalu jabarkanlah sepanjang baris terakhirnya
(dengan tanda $(-1)^{n+n} = +1$): maka determinan sisanya adalah
$C_{n-1}$. Dengan menghimpun faktor pertanyaan 11--12:
\[
C_n = \frac{\prod_{i<n}(a_n - a_i)\,\prod_{j<n}(b_n - b_j)}
{\prod_{j}(a_n + b_j)\,\prod_{i<n}(a_i + b_n)}\;C_{n-1},
\]
dan induksinya (dengan \omterm{def:b1:vspaces:free}{basis} $C_1 = \frac1{a_1+b_1}$) merakit
tepat alternan ganda Cauchy: yaitu faktor $(a_j - a_i)(b_j
- b_i)$ bagi semua pasangan, dibagi semua jumlah $(a_i + b_j)$.

\textbf{13.} Rumusnya lenyap jika dan hanya jika ada $a_j = a_i$ atau $b_j =
b_i$: sehingga matriks Cauchy dapat dibalik jika dan hanya jika kedua keluarganya
berbeda berpasangan. Hilbert: $a_i = i$, $b_j = j - 1$. Untuk $n =
2$: pembilangnya $(2-1)(1-0) = 1$, penyebutnya $1\cdot2\cdot2\cdot3
= 12$: jadi $\det H_2 = \frac1{12}$. Untuk $n = 3$: pembilangnya
$\bigl[(1)(2)(1)\bigr]^2 = 4$, penyebutnya $(1\cdot2\cdot3)
(2\cdot3\cdot4)(3\cdot4\cdot5) = 6\cdot24\cdot60 = 8640$:
jadi $\det H_3 = \frac{4}{8640} = \frac1{2160}$. Inversnya untuk $n = 2$:
\[
H_2^{-1} = 12\begin{pmatrix} \frac13 & -\frac12\\[2pt]
-\frac12 & 1\end{pmatrix}
= \begin{pmatrix} 4 & -6\\ -6 & 12 \end{pmatrix},
\]
semuanya bilangan bulat (yaitu gejala yang benar untuk setiap $n$).

\textbf{14.} Matriks sistemnya adalah matriks Cauchy,
yang dapat dibalik menurut pertanyaan 13 ketika $b_j$ (dan $a_i$)
berbeda berpasangan: jadi penyelesaiannya tunggal. Tafsirannya: sebuah fungsi
rasional $R = \sum_j \frac{c_j}{X + b_j}$ berkutub sederhana
ditentukan oleh $n$ nilainya $R(a_1), \dots, R(a_n)$, dan
sebaliknya setiap lembar data yang demikian terwujud tepat sekali ---
yaitu padanan pencuplikan atas teorema keberadaan dan ketunggalan
pecahan parsialnya (\cref{thm:b1:fractions:complex}).

\textbf{15.} Vieta bagi $X^3 + pX + q$: $\lambda_1 + \lambda_2 +
\lambda_3 = 0$, $\sum_{i<j}\lambda_i\lambda_j = p$, sehingga $p_1 = 0$
dan $p_2 = p_1^2 - 2p = -2p$. Setiap akarnya memenuhi $\lambda^3 =
-p\lambda - q$; dan menjumlahkannya: $p_3 = -p\,p_1 - 3q = -3q$. Mengalikannya
dengan $\lambda$ lalu menjumlahkannya: $p_4 = -p\,p_2 - q\,p_1 = 2p^2$.

\textbf{16.} Menurut pertanyaan 5 (karena kesamaan $S = W^{\mathsf T}W$
dan $\det S = V^2$ sah atas $\C$),
\[
V^2 = \begin{vmatrix}
3 & 0 & -2p\\
0 & -2p & -3q\\
-2p & -3q & 2p^2
\end{vmatrix}
= 3\bigl(-4p^3 - 9q^2\bigr) + (-2p)\bigl(0 - 4p^2\bigr)
= -4p^3 - 27q^2 ,
\]
dengan menjabarkannya sepanjang baris pertamanya.

\textbf{17.} Akar rangkap berarti dua $\lambda_i$ yang sama, yaitu
$V = 0$, yaitu $\operatorname{disc} = -4p^3 - 27q^2 = 0$. Bagi
$X^3 - 3X + 2$: $4(-3)^3 + 27\cdot4 = -108 + 108 = 0$, yang cocok dengan
akar rangkap $1$ pada $(X-1)^2(X+2)$.

\textbf{18.} Akar yang tak real pada kubik real datang berpasangan
sekawan, sehingga tepat dua kasus terjadi ketika $\operatorname{disc} \neq
0$. Tiga akar real yang berbeda: maka $V$ real dan taknol, sehingga
$\operatorname{disc} = V^2 > 0$. Satu akar real $\lambda_1$ dan
$\lambda_3 = \conj{\lambda_2} \notin \R$: maka
\[
(\lambda_2 - \lambda_1)(\lambda_3 - \lambda_1) =
\abs{\lambda_2 - \lambda_1}^2 > 0,
\qquad
\lambda_3 - \lambda_2 = -2\iu\,\operatorname{Im}\lambda_2 \neq 0,
\]
sehingga $V$ bilangan murni khayal yang taknol dan
$\operatorname{disc} = V^2 < 0$. Jadi kedua tandanya mencirikan
kedua kasusnya.

\textbf{19.} Ambillah $b_i = a_i$ pada alternan gandanya: maka
pembilangnya $\prod_{i<j}(a_j - a_i)^2 > 0$ dan penyebutnya
$\prod_{i,j}(a_i + a_j) > 0$ (karena semua entrinya positif): sehingga
determinannya positif. (Dalam bahasa kemudian: kernel
$\frac1{x+y}$ bersifat tegas positif.)

\textbf{20.} Menurut pertanyaan 2--3, determinannya sama dengan
$V(2,4,7)/(0!\,1!\,2!) = \frac{(4-2)(7-2)(7-4)}{2} =
\frac{30}{2} = 15$. Secara langsung, matriksnya adalah
\[
\begin{pmatrix}
1 & 2 & 1\\
1 & 4 & 6\\
1 & 7 & 21
\end{pmatrix},
\qquad
\det = (84 - 42) - 2(21 - 6) + (7 - 4) = 42 - 30 + 3 = 15 .
\]

\textbf{21.} Andaikan $\sum_i c_i\,(\lambda_i^{\,k})_{k} = 0$ sebagai
sebuah barisan. Membacanya di $k = 0, 1, \dots, n-1$ memberikan $W^{\mathsf
T}c = 0$ dengan $W = (\lambda_i^{\,j-1})$ yang dapat dibalik ($\det = V
\neq 0$, karena $\lambda_i$ berbeda): sehingga $c = 0$. Jadi barisan geometrinya
\omterm{def:b1:vspaces:free}{bebas}.

\textbf{22.} $a = b = (1, 2, 3)$: pembilangnya $\bigl[(2-1)(3-1)
(3-2)\bigr]^2 = 4$; penyebutnya $\prod_{i,j}(i + j) =
(2\cdot3\cdot4)(3\cdot4\cdot5)(4\cdot5\cdot6) = 24\cdot60\cdot120
= 172800$. Maka $\det\bigl(\frac1{i+j}\bigr) = \frac{4}{172800}
= \frac1{43200}$.

\textbf{23.} Jika $x_i = x_j$, maka penukaran kedua variabelnya menetapkan
titiknya tetapi mesti mengubah tanda $F$: jadi $F = -F$, sehingga $F = 0$
di situ. \omterm{def:b1:arith:divides}{Keterbagiannya}: pandanglah $F$ sebagai \omterm{def:b1:poly:def}{polinomial} pada satu
variabel $x_n$ dengan koefisien pada variabel lainnya; ia
lenyap pada $n - 1$ ``nilai'' $x_1, \dots, x_{n-1}$, sehingga
pemfaktoran berulang (\cref{thm:b1:poly:factor}) memberikan $F =
\prod_{i<n}(x_n - x_i)\cdot G$ dengan $G$ \omterm{def:b1:poly:def}{polinomial}. Adapun
prafaktornya invarian terhadap penukaran dua indeks $i, j < n$, sehingga
$G$ berselang-seling pada $x_1, \dots, x_{n-1}$, dan induksinya
melengkapkannya: yaitu $\prod_{i<j}(x_j - x_i)$ \omterm{def:b1:arith:divides}{membagi} $F$.

\textbf{24.} $D = \det(x_i^{\,j-1})$ merupakan \omterm{def:b1:poly:def}{polinomial} pada
$x_i$; dan menukar dua variabelnya menukar dua barisnya, sehingga $D$
berselang-seling, jadi menurut pertanyaan 23, $D = c\,\prod_{i<j}(x_j - x_i)$
bagi suatu \omterm{def:b1:poly:def}{polinomial} $c$. Derajat totalnya: $D$ berderajat $\leq 0 +
1 + \dots + (n-1) = \binom n2$, sedangkan hasil kalinya berderajat tepat
$\binom n2$: sehingga $c$ sebuah tetapan. Adapun monomial $x_2\,x_3^2\cdots
x_n^{\,n-1}$ berkoefisien $1$ pada $D$ (yaitu hasil kali diagonalnya) dan
$1$ pada hasil kalinya (pilihlah variabel berindeks lebih besar pada setiap
faktornya): jadi $c = 1$, dan rumus Vandermondenya jatuh tanpa
induksi.

\textbf{25.} (i) Sifat berselang-selingnya --- yaitu aksioma ``dua kolom yang sama
membunuh determinannya'' --- itulah mesinnya: karena ia menghasilkan setiap
faktor $(x_j - x_i)$, $(a_j - a_i)$, $(b_j - b_i)$ pada
soal ini. (ii) Kesamaan $\det S = V^2$ mengganti
akarnya yang kompleks dan tak terjangkau satu per satu dengan jumlah pangkatnya,
yang berupa \omterm{def:b1:poly:def}{polinomial} real pada koefisiennya, sehingga keberbedaannya
menjadi tanda sebuah bilangan real yang dapat dihitung. (iii) Adapun
\omterm{ex:b1:det:vandermonde}{determinan Vandermonde} menguasai interpolasi \omterm{def:b1:poly:def}{polinomial}, sedangkan
\omterm{pb:b1:det:1}{determinan Cauchy} menguasai pecahan parsial dan fungsi rasional
yang tercuplik (dengan \omterm{pb:b1:det:1}{matriks Hilbert} sebagai kasus khususnya yang
paling termasyhur). (iv) Setiap \omterm{def:b1:poly:def}{polinomial} berselang-seling terbagi oleh
$\prod_{i<j}(x_j - x_i)$, lalu cacah derajatnya memakukan \omterm{def:b1:poly:def}{polinomial}
yang demikian sampai sebuah tetapan --- dan itulah sebabnya hasil kali ini
terus muncul kembali di mana pun sebuah determinan lenyap pada
keberimpitan. Adapun teorema Bagian III adalah \emph{alternan ganda
Cauchy}.

\end{solution}
